% ExampleHyouka.tex

\documentclass[a3paper, twocolumn]{article}
\input{ExperimentHeader.tex}

% 以下内容全部由ai生成，仅用以演示

\pgfkeys{/reportinformation/title = 关于电影剧本结局的推理委托实验}

\begin{document}

\maketitle

\section{实验目的}
\begin{itemize}
    \item 检验仅凭剧本中既存线索进行逻辑推理能否准确还原原作者（本乡）预设的结局；
    \item 探究委托人（入须冬实）的真实意图对推理过程及结论的影响；
    \item 观察推理者（折木奉太郎）在不知情情况下代行创作决策时产生的认知偏差。
\end{itemize}

\section{实验准备工作}
\subsection{实验材料}
\begin{itemize}
    \item 电影《万人的死角》未完成剧本（前中期剧情完整，结局缺失）；
    \item 剧本原作者本乡的写作习惯及前期情节设定摘要；
    \item 已完成拍摄的影片片段及分镜稿；
    \item 知情人江波（本乡好友）的存在信息（初始未公开）。
\end{itemize}

\subsection{实验人员}
\begin{itemize}
    \item 委托人：入须冬实（班级文化祭实行委员，负责统筹电影制作）；
    \item 推理者：折木奉太郎（古典部成员，擅长逻辑推理）；
    \item 观察员：千反田爱瑠、福部里志、伊原摩耶花（古典部其他成员）；
    \item 隐性知情人：江波（知晓本乡原定结局，但初期未被引入实验）。
\end{itemize}

\subsection{委托陈述}
入须冬实向折木奉太郎声称：原作者本乡因故无法写出结局，需要折木从已有线索中“推理”出本乡原本设定的真凶及结局，以指导后续拍摄。

\section{实验过程}
\subsection{阶段一：线索收集与逻辑推演}
折木奉太郎观看已完成片段，梳理剧本中的伏笔、人物动机及时间线。基于本乡前期偏好的“意外但合理”的风格，排除几个明显矛盾的嫌疑人，最终推演出一个逻辑自洽的“真凶”——将凶手指向一名看似无关的配角。

\subsection{阶段二：结论交付与拍摄推进}
折木将推理结果告知入须。入须表示满意，并迅速安排剧组按照折木的结局进行拍摄。此时古典部其他成员均认为委托已成功解决。

\subsection{阶段三：关键质疑与信息缺口暴露}
千反田爱瑠在事后提出一个核心疑问：既然本乡的好友江波一直参与剧本讨论，且江波本人明确表示“知道本乡原本的结局”，为什么入须不直接询问江波，反而要绕远路请折木“推理”？这一逻辑矛盾引起折木的重新审视。

\subsection{阶段四：真相揭示与认知重构}
折木通过追问江波及再次分析入须的行为模式，发现真相：
\begin{itemize}
    \item 本乡的原定结局“平淡而写实”（例如真凶是意外死亡，无人承担罪责）；
    \item 入须认为这样的结局“缺乏戏剧性”，无法在文化祭获得好评；
    \item 入须的真实目的并非“还原本乡原意”，而是希望得到一个“更精彩、更适合拍成电影的结局”；
    \item 她利用折木“擅长推理”的名声，让折木在不知情的情况下代行编剧职能，从而绕过“擅自篡改剧本”的道德争议。
\end{itemize}
折木的推理结论并非本乡的原定结局，而是入须期望的“高完成度悬疑结局”。

\section{实验结果}
\subsection{主要发现}
\begin{enumerate}
    \item 折木奉太郎成功完成了一次逻辑自洽的推理，但该结论与原作者本乡的真实意图完全无关。
    \item 实验揭示委托人入须冬实的真实目的是“借用推理者的能力进行剧本创作”，而非“求解原设答案”。
    \item 推理者折木在得知自己被利用后，产生了强烈的被背叛感，并一度对人际关系产生不信任。
    \item 入须冬实承认了利用行为，并辩解称这是为了集体利益（班级电影顺利完成并获得好评）而做出的必要策略。
\end{enumerate}

\subsection{实验结论}
本实验表明：在“推理委托”中，委托人的动机预设会极大影响推理的有效方向。若委托表面是“还原真相”，实质是“要求创作”，则推理者即使逻辑无误，其结论仍可能完全偏离事实原点。这种认知错位可称为“伪委托陷阱”。


\end{document}

%%% Local Variables:
%%% mode: LaTeX
%%% TeX-master: t
%%% End:
